% Final presentation, CSE 8803 NBD (Fall 2026).
% Build: `make` in this directory -> slides.pdf (classroom copy, no notes), presentation_andra-thomas.pdf (identical submission copy)
% and slides_notes.pdf (printable handout: each slide with its speaker notes below).
%
% Number sources (every number on a slide comes from one of these; paths relative to the repo root):
%   RESULTS.md                                  summary tables (primary, secondary, robustness, pilots A12-A14)
%   results/gate/summary.json                   prior gate (full population)
%   results/stats/d5/{primary,secondary,matrix,drift}.csv     primary analysis, d = 5
%   results_residual/stats/d5/*.csv             trial-locked residual rerun
%   results/stats/d{3,10}/primary.csv           d sensitivity
%   results{,_residual}/robustness/d5.csv       leave-one-mouse-out, drop DY_014 (post hoc)
%   results/stats/d5/geometry_sessions.csv      per-session angles (animal means used on the geometry slide)
%   results/subspaces/*_d5.json                 covariance-matched chance angle (floor_mean_deg_mean, session mean)
%   results/estimator/*_d5.pqt                  estimator check (session means of r2_session_alpha, as in RESULTS.md)
%   results/pilot/summary_d5.json, sessions_d5.csv      A12 (stable units, probe motion, pupil)
%   results/a13/summary_d5.json                 A13 (engagement, satiety, amplitude, multi-unit)
%   results/a14/summary_d5.json                 A14 (amplitude coupling, corrected refits, timing)
%   results_{stable,mua,ampstable,ampcorr_low,ampcorr_high,ampcorr_slow}/stats/d5/{matrix,drift}.csv
%   results/a15/summary_d5.json, results_region2/stats/d5/*.csv, data_region2/meta/frozen.json   A15 second region
%   figures/slide_fig5.pdf, figures/slide_figS1.pdf    made by scripts/20_slide_figures.py from the files above
%   The controls chart (pgfplots) adds same-session originals: results/stats/d5/drift.csv restricted to each
%   version's sessions, animal means (rate-stable 0.834, multi-unit 0.879, amplitude-stable 0.734).
\documentclass[aspectratio=169,11pt]{beamer}

\usetheme[progressbar=frametitle,numbering=fraction,block=fill,sectionpage=none]{metropolis}
\usepackage{appendixnumberbeamer}
\usepackage[T1]{fontenc}
\usepackage[sfdefault,lining]{FiraSans}
\usepackage{newtxsf}
\usepackage{booktabs}
\usepackage{array}
\usepackage{tikz}
\usetikzlibrary{arrows.meta,positioning,calc,decorations.pathreplacing}
\usepackage{pgfplots}
\pgfplotsset{compat=1.18}
\usepackage{graphicx}

% ---------------------------------------------------------------- notes
\ifdefined\withnotes
  % Presenter copy for printing or reading on a phone: each portrait page = slide on top, its notes below.
  \usepackage{pgfpages}
  \setbeameroption{show notes}
  \pgfpagesuselayout{2 on 1}[letterpaper, border shrink=6mm]
\fi
\setbeamertemplate{note page}{%
  \vspace*{2mm}\noindent\hspace*{4mm}%
  \begin{minipage}[t]{\dimexpr\paperwidth-24mm\relax}
    \raggedright
    {\color{mDarkTeal}\bfseries\insertframetitle\par}
    \vspace{1mm}
    \small\insertnote
  \end{minipage}}
\newcommand{\why}[1]{\par\smallskip{\footnotesize\textbf{Why these choices (for you, not to read out):} #1\par}}
\newcommand{\qa}[2]{\par\smallskip\textbf{Likely question:} #1\par\textbf{Answer:} #2}

% ---------------------------------------------------------------- colours (dataviz reference slots 1-4, as in figures/)
\definecolor{cRRR}{HTML}{2A78D6}
\definecolor{cPCA}{HTML}{EB6834}
\definecolor{cRand}{HTML}{1BAF7A}
\definecolor{cRidge}{HTML}{EDA100}
\definecolor{cGrey}{HTML}{9A9893}
\definecolor{cInk}{HTML}{52514E}
\setbeamercolor{block title}{bg=mDarkTeal!12,fg=mDarkTeal}
% Slightly slimmer frame-title bar, and a title page that fits the plain frame (metropolis uses \paperheight,
% which always overflows the frame by about 16 pt).
\makeatletter
\setlength{\metropolis@frametitle@padding}{1.6ex}
\setbeamertemplate{title page}{%
  \begin{minipage}[b][\dimexpr\paperheight-18pt\relax]{\textwidth}
    \vfill%
    \usebeamertemplate*{title}\usebeamertemplate*{subtitle}\usebeamertemplate*{title separator}%
    \usebeamertemplate*{author}\usebeamertemplate*{date}\usebeamertemplate*{institute}%
    \vfill
  \end{minipage}}
\makeatother
\setbeamerfont{institute}{size=\small}
\setbeamercolor{block body}{bg=mDarkTeal!4}

\newcommand{\hl}[1]{{\color{mDarkTeal}\bfseries #1}}
\newcommand{\muted}[1]{{\color{cInk}#1}}
\newcommand{\Rsq}{$R^2$}
\newcolumntype{L}[1]{>{\raggedright\arraybackslash}p{#1}}

% Figure with a placeholder fallback so the deck compiles before the figure exists.
% #1 path, #2 width, #3 height
\newcommand{\figorplaceholder}[3]{%
  \IfFileExists{#1}{\includegraphics[width=#2,height=#3,keepaspectratio]{#1}}{%
    \setlength{\fboxsep}{0pt}\fcolorbox{cGrey}{white}{\parbox[c][#3][c]{#2}{\centering\small\color{cInk}
      Figure pending:\\[2pt]\texttt{\detokenize{#1}}}}}}

\pgfplotsset{
  slide axis/.style={
    font=\small,
    axis lines*=left,
    tick style={draw=none},
    ymajorgrids, grid style={cGrey!35},
    axis line style={cGrey},
    every tick label/.append style={font=\small, color=cInk},
    label style={font=\small, color=cInk},
    enlarge x limits=0.1,
  },
  ci bars/.style={error bars/.cd, y dir=both, y explicit, error bar style={line width=0.8pt, black!70},
    error mark options={rotate=90, mark size=2.5pt, black!70}},
  subspace bars/.style={slide axis, ybar, bar width=6.5mm, width=0.94\linewidth, height=5.4cm,
    symbolic x coords={r1,r4,r16,r60,pca,rnd}, xtick={r1,r4,r16,r60,pca,rnd},
    xticklabels={{RRR\\0.5 s},{RRR\\2 s},{RRR\\8 s},{RRR\\30 s},PCA,random},
    xticklabel style={align=center, text depth=0pt}, yticklabel style={/pgf/number format/fixed},
    enlarge x limits=0.07},
}

\title{Prediction Horizon and the Content of Neural Population Subspaces}
\subtitle{What long prediction horizons find in mouse secondary motor cortex}
\author{Gunnar Andra-Thomas}
\institute{CSE 8803: Statistical Machine Learning Models for Neural and Behavioral Data Analysis\\
Georgia Tech \textperiodcentered{} Final project presentation\\[4pt]
{\color{mDarkTeal}Introduction $\rightarrow$ Method $\rightarrow$ Results and interpretation $\rightarrow$ Conclusion}}
\date{Monday 23 November 2026}
\hypersetup{pdftitle={Prediction Horizon and the Content of Neural Population Subspaces},
  pdfauthor={Gunnar Andra-Thomas}, pdfsubject={CSE 8803 final project presentation}}

\begin{document}

% =====================================================================================
\maketitle
\note{%
Hi, I'm Gunnar. My project asks a simple question: if you look for the directions in a neural population that best
predict its own future, does how far ahead you predict decide what those directions encode? In mouse secondary motor
cortex the answer is mostly no. The long horizon finds a slow mode that tracks time in the session, and the second half
of the talk is about what that mode is and is not. I'll set up the question and the method, go through the three
pre-registered tests, then ask whether the slow mode is a recording artifact and whether it generalises.}

% =====================================================================================
\section{Introduction}
\begin{frame}{Why ask this?}
\begin{columns}[T,onlytextwidth]
\column{0.48\textwidth}
\textbf{Neuroscience}

\medskip
Animals carry beliefs about hidden task states across many trials. Is such a belief a \hl{slow dynamical mode}
that the circuit carries forward in time, or a \hl{static code}?

\medskip
The IBL block prior is decodable across the brain (Findling et al.\ 2025): a good test case.

\column{0.48\textwidth}
\textbf{Methods}

\medskip
Predictive objectives (lagged regression, slow-feature and dynamical-components analysis, predictive learning in ML)
find structure \hl{without labels}.

\medskip
If the horizon selected content, it would be a cheap way to find latent-state subspaces. If not, we should know what
long horizons \emph{do} find.
\end{columns}
\note{%
Why ask this? Animals keep beliefs about hidden task states over many trials, and a basic question is how: is the
belief a slow dynamical mode that the circuit carries forward, or a static code? The IBL block prior is a good test
case, because it is decodable across the brain. On the methods side, predictive objectives are a common label-free way
to find structure. If the prediction horizon decided content, that would be a cheap way to find latent-state
subspaces; if not, we should know what long horizons actually find.}
\end{frame}

% =====================================================================================
\begin{frame}{Does the prediction horizon decide what a subspace encodes?}
\begin{columns}[T,onlytextwidth]
\column{0.46\textwidth}
\textbf{IBL decision task.} Mice report the side of a grating. The side is drawn from unsignalled blocks with
$P(\text{left}) = 0.8$ or $0.2$.

\medskip
\begin{tikzpicture}
\begin{axis}[slide axis, width=\linewidth, height=3.7cm, xmin=0, xmax=150, ymin=0, ymax=1,
  ytick={0.2,0.8}, xtick=\empty, xlabel={trials}, ylabel={$P(\text{left})$}, ymajorgrids=false,
  enlarge x limits=0, ylabel style={yshift=-4pt}]
\addplot[const plot, cGrey, line width=1.2pt] coordinates {(0,0.8) (45,0.2) (100,0.8) (150,0.8)};
\addplot[cRRR, line width=1.4pt, smooth] coordinates {(0,0.533) (6,0.668) (12,0.735) (18,0.768) (24,0.784) (30,0.792) (36,0.796) (42,0.798) (45,0.732) (48,0.574) (51,0.462) (54,0.384) (60,0.291) (66,0.245) (72,0.222) (78,0.211) (84,0.205) (90,0.203) (99,0.201) (102,0.379) (105,0.504) (108,0.592) (114,0.698) (120,0.750) (126,0.775) (132,0.788) (138,0.794) (147,0.798)};
\node[font=\small, color=cInk, anchor=south west] at (axis cs:2,0.82) {block};
\node[font=\small, color=cRRR, anchor=west, align=left] at (axis cs:106,0.3) {Bayes\\prior};
\end{axis}
\end{tikzpicture}

\muted{\small Schematic. The prior is a slow latent the mouse must infer from trial history.}

\column{0.5\textwidth}
\begin{block}{Hypothesis}
Predicting \hl{30 s} ahead isolates a slow subspace carrying the \hl{block prior}; predicting \hl{0.5 s} ahead
isolates a fast one carrying the \hl{current stimulus}.
\end{block}
\medskip
\textbf{Pre-registered tests}\\[2pt]
\hl{P1} 0.5 s and 30 s subspaces differ\\
\hl{P2} prior decoding rises from 0.5 s to 30 s\\
\hl{P3} stimulus decoding falls from 0.5 s to 30 s
\end{columns}
\note{%
The task is the IBL visual decision task. The stimulus side is biased 80/20 in blocks that switch without any cue, so
the mouse has to infer the block. The Bayes-optimal prior, in blue, is that inference: a slow variable that lags the
true block. The hypothesis was that a long prediction horizon pulls out this kind of slow latent, and a short horizon
pulls out fast sensorimotor content. That became three tests, fixed before fitting any real data: the subspaces
differ, prior decoding rises with horizon, and stimulus decoding falls. The second and third together are the core
of the hypothesis: content should sort by timescale.}
\end{frame}

% =====================================================================================
\section{Method}
\begin{frame}{Method: predict activity $\Delta$ seconds ahead through a rank-$d$ bottleneck}
\small\textbf{Data:} IBL Brain-Wide Map \textperiodcentered{} MOs \textperiodcentered{} 16 sessions / 13 mice
\textperiodcentered{} 22--128 units \textperiodcentered{} 0.5 s bins

\centering
\begin{tikzpicture}[x=1cm,y=1cm,
  unit/.style={circle, draw=cInk, fill=white, minimum size=3.0mm, inner sep=0pt},
  lat/.style={rectangle, rounded corners=1pt, draw=cRRR, fill=cRRR!20, minimum size=3.0mm, inner sep=0pt},
  conn/.style={cGrey!70, line width=0.35pt}]
  \draw[-{Stealth[length=2mm]}, cInk, line width=0.8pt] (-0.5,1.95) -- (6.9,1.95) node[right, font=\small, color=cInk] {time};
  \draw[cInk, line width=0.8pt] (0,1.82) -- (0,2.08); \node[font=\small, above] at (0,2.0) {$t$};
  \draw[cInk, line width=0.8pt] (6,1.82) -- (6,2.08); \node[font=\small, above] at (6,2.0) {$t+\Delta$};
  \draw[decorate, decoration={brace, amplitude=3pt, mirror}, cInk] (0.15,1.78) -- (5.85,1.78)
     node[midway, below=2pt, font=\small, color=cInk] {$\Delta \in \{0.5,\ 2,\ 8,\ 30\}$ s};
  \foreach \i in {0,...,5} {
    \node[unit] (a\i) at (0, 1.25-0.24*\i) {};
    \node[unit] (c\i) at (6, 1.25-0.24*\i) {};
  }
  \foreach \j in {0,...,2} { \node[lat] (b\j) at (3, 0.89-0.24*\j) {}; }
  \foreach \i in {0,...,5} { \foreach \j in {0,...,2} {
      \draw[conn] (a\i) -- (b\j); \draw[conn] (b\j) -- (c\i); } }
  \node[font=\small, below=1mm of a5] {$x_t$ now};
  \node[font=\small, below=1mm of c5] {$\hat x_{t+\Delta}$};
  \node[font=\small, color=cRRR, below=5.5mm of b2] {$z_t = V_\Delta^\top x_t$, $d = 5$};
  \node[font=\small, color=cInk] at (1.35,0.0) {$V_\Delta$};
  \node[font=\small, color=cInk] at (4.65,0.0) {$W_\Delta$};
  \node[anchor=north west, text width=5.85cm, align=left, font=\normalsize] at (7.25,1.75) {%
    The \hl{subspace} is the span of $V_\Delta$: the directions of current activity that best predict activity
    $\Delta$ later.\\[5pt]
    Then decode task variables from $z_t$.};
\end{tikzpicture}

$\hat B_\Delta=\arg\min\nolimits_{\operatorname{rank}(B)\le d}\ \lVert X_{t+\Delta}-X_tB\rVert_F^2+\lambda\lVert B\rVert_F^2,
\qquad \hat B_\Delta = V_\Delta W_\Delta^\top$

\raggedright\muted{\small Tests committed before any real-data fit; later changes logged as amendments. Null: 200
pseudo-sessions. Stats: one-sided Wilcoxon across 13 mice, corrected for 3 tests (Holm).}
\note{%
The method is lagged reduced-rank regression, a form of supervised dimensionality reduction with PCA as the
unsupervised baseline. MOs was the only candidate region with enough sessions.
Squeeze the population now through five latents and use them to predict the population
$\Delta$ seconds later, at 0.5, 2, 8 and 30 seconds. $V_\Delta$ defines the subspace; then I decode task variables
from the latents. Three safeguards against slow signals posing as cognition: tests fixed in advance; every
experimenter-defined variable scored against pseudo-sessions, block sequences re-drawn from the task's own generative
process; and statistics across mice, not sessions.
\why{RRR is closed form (ridge fit, then the top $d$ singular directions of the fitted predictions): no local optima,
only $d$ and $\lambda$. One cross-validated $\lambda$ per session is shared by all lags, because $\lambda$ itself moves
the subspace (CCA-like when small, PLS-like when large). Pseudo-session null: within a session any slow signal
correlates with slowly switching blocks by chance (Harris 2020). Wilcoxon across mice: sessions from one mouse are not
independent; a rank test assumes little with $n = 13$. Holm: family-wise control over the 3 tests, less conservative
than Bonferroni.}}
\end{frame}

% =====================================================================================
\section{Results and interpretation}
\begin{frame}{Sanity checks: the prior is in MOs, and the estimator works}
\begin{columns}[T,onlytextwidth]
\column{0.36\textwidth}
\textbf{Prior, full population}\\
\muted{\small Findling-style decoder vs.\ pseudo-session null}

\medskip
{\Large\hl{0.064}}\\ median corrected \Rsq{}\\[4pt]
{\Large\hl{5 of 16}}\\ sessions significant alone\\[4pt]
{\Large\hl{$1.2\times10^{-5}$}}\\ Fisher-combined $p$

\medskip
Replicates Findling et al.\ 2025.

\column{0.60\textwidth}
\centering
\begin{tikzpicture}
\begin{axis}[slide axis, width=0.8\linewidth, height=5.0cm, xmode=log, log basis x=2,
  xtick={0.5,2,8,30}, xticklabels={0.5 s,2 s,8 s,30 s}, xmin=0.4, xmax=220, enlarge x limits=0, clip=false,
  ymin=0, ymax=0.4, ytick={0,0.1,0.2,0.3,0.4}, yticklabel style={/pgf/number format/fixed},
  xlabel={prediction lag $\Delta$}, ylabel={held-out \Rsq{}, future activity}]
% session means of r2_session_alpha, results/estimator/*_d5.pqt (= RESULTS.md estimator table at 0.5 s and 30 s)
\addplot[cRidge, line width=1.4pt, mark=*, mark size=2pt] coordinates {(0.5,0.371) (2,0.225) (8,0.202) (30,0.178)};
\addplot[cRRR, line width=1.4pt, mark=*, mark size=2pt] coordinates {(0.5,0.246) (2,0.194) (8,0.185) (30,0.176)};
\addplot[cPCA, line width=1.4pt, mark=*, mark size=2pt] coordinates {(0.5,0.232) (2,0.172) (8,0.162) (30,0.152)};
\addplot[cRand, line width=1.4pt, mark=*, mark size=2pt] coordinates {(0.5,0.163) (2,0.135) (8,0.123) (30,0.121)};
\node[font=\small, anchor=west, color=cInk] at (axis cs:0.62,0.372) {full-rank ridge};
\node[font=\small, anchor=west, color=cInk] at (axis cs:34,0.186) {RRR, $d = 5$};
\node[font=\small, anchor=west, color=cInk] at (axis cs:34,0.150) {PCA, $d = 5$};
\node[font=\small, anchor=west, color=cInk] at (axis cs:34,0.115) {random, $d = 5$};
\end{axis}
\end{tikzpicture}

\raggedright
RRR beats PCA at every lag (all $p < 0.01$). At 30 s, \hl{5 dimensions recover 96\%} of full-rank ridge.
\end{columns}
\note{%
Two sanity checks. First, the prior really is in MOs: decoding it from the whole population, as Findling and colleagues
did, beats the pseudo-session null, combined $p$ around $10^{-5}$, though the effect is small. Second, the estimator
works: lagged subspaces predict the future better than PCA or random subspaces of the same size, and at 30 seconds five
dimensions reach 96\% of an unconstrained model. So long-horizon predictable structure is very low-dimensional, and the rank limit is not what hides anything that
is predictable. One
note on process: after the first run I found a normalisation bug; it was fixed and logged, and the primary results did
not change.
\qa{Why RRR rather than PCA?}{It makes the horizon an explicit knob with everything else fixed. Its edge over PCA is
modest and mostly trial-locked structure (backup).}}
\end{frame}

% =====================================================================================
\begin{frame}{Geometry (P1): the subspaces differ, mostly via the short lag}
\begin{columns}[T,onlytextwidth]
\column{0.48\textwidth}
\centering
\begin{tikzpicture}
\begin{axis}[slide axis, ybar, bar width=13mm, width=\linewidth, height=5.6cm,
  symbolic x coords={s05,cross,s30}, xtick={s05,cross,s30},
  xticklabels={{same lag\\0.5 s},{cross lag\\0.5 vs 30 s},{same lag\\30 s}},
  xticklabel style={align=center}, enlarge x limits=0.25,
  ymin=0, ymax=90, ytick={0,30,60,90}, ylabel={mean principal angle (deg)}]
% animal means over 13 mice, results/stats/d5/geometry_sessions.csv
\addplot[fill=cRRR, draw=none] coordinates {(s05,58.9) (cross,71.6) (s30,68.0)};
% covariance-matched chance: session mean of floor_mean_deg_mean, results/subspaces/*_d5.json
\draw[cInk, dashed, line width=1pt] (rel axis cs:0,0.7533) -- (rel axis cs:1,0.7533);
\node[font=\small, color=cInk, anchor=south west] at (rel axis cs:0.0,0.80) {chance 67.8\textdegree};
\end{axis}
\end{tikzpicture}

\column{0.48\textwidth}
Fit each lag on \hl{opposite halves} of a session. \hl{Floor} = angle between the two halves' fits of the same lag.

\medskip
Cross-lag beats both floors (both required):\\
Holm $p = 0.0026$; 11 of 13 mice.\\[3pt]
\quad vs.\ 0.5 s floor: $+12.7^\circ$ [10.2, 15.1]\\
\quad vs.\ 30 s floor: $+3.6^\circ$ [2.2, 4.9]

\medskip
\muted{The 30 s floor is itself at chance; the effect fails in the predicted-activity basis ($p = 0.96$).}
\end{columns}
\note{%
First test, geometry: fit each subspace on one half of the session and compare with fits on the other half. Short
lags are reproducible, about 59 degrees, below the 68-degree chance; 30 s subspaces sit at chance. The cross-lag angle beats both floors, so the test passes, but mostly through the short lag; scientifically it is
the weakest of the three results.
The leading slow direction is shared across lags; the other four dimensions differ as much as random ones, and the
30 s ones are not even reproducible between halves.
\why{Chance: random subspaces drawn in proportion to the data covariance (Elsayed and Cunningham 2017).
``Both required'' is an intersection--union test: $p$ = the larger of the two one-sided $p$-values.}
\qa{Does it depend on the basis?}{Yes. In the loading basis (predicted activity) P1 fails, $p = 0.96$,
with cross-lag minus the 30 s floor $-1.0^\circ$.}
\qa{Why is the 30 s floor at chance?}{About 120 non-overlapping 30 s windows per session, and halves confound
reliability with nonstationarity.}}
\end{frame}

% =====================================================================================
\begin{frame}{Content (P3): long-lag subspaces lose the current stimulus}
\begin{columns}[T,onlytextwidth]
\column{0.55\textwidth}
\centering
\begin{tikzpicture}
\begin{axis}[subspace bars, ymin=0, ymax=0.32, ytick={0,0.1,0.2,0.3}, ylabel={stimulus side, corrected AUC}]
% results/stats/d5/matrix.csv, variable cur_side (mean and 95% CI over 13 animals)
\addplot[fill=cRRR, draw=none, bar shift=0pt, ci bars] coordinates {
  (r1,0.218) += (0,0.050) -= (0,0.049)
  (r4,0.161) += (0,0.032) -= (0,0.031)
  (r16,0.144) += (0,0.042) -= (0,0.036)
  (r60,0.158) += (0,0.052) -= (0,0.048)};
\addplot[fill=cGrey, draw=none, bar shift=0pt, ci bars] coordinates {
  (pca,0.240) += (0,0.057) -= (0,0.059)
  (rnd,0.218) += (0,0.058) -= (0,0.053)};
\end{axis}
\end{tikzpicture}

\muted{\small Corrected = AUC minus pseudo-session null median; 95\% CI over mice}

\column{0.41\textwidth}
Falls by \hl{0.061} [0.028, 0.094] from 0.5 s to 30 s: Holm $p = 0.0061$; 10 of 13 mice.

\medskip
Holds on trial-locked residuals ($p = 0.0051$) and at $d = 10$; not at $d = 3$ ($p = 0.11$).

\medskip
\muted{The 0.5 s subspace is not \emph{specialised}: it equals a random subspace (0.218). Long lags lose the stimulus;
short lags do not gain it.}
\end{columns}
\note{%
Now content. For each subspace I decode a variable from its five latents and subtract what the same decoder gets on
pseudo-sessions, so zero means no information beyond chance. The current stimulus side falls as the horizon grows,
significantly across mice. The nuance: the 0.5 second subspace is no better than a random five-dimensional one; it is
the long lags that lose the stimulus. That is roughly what evoked responses lasting under a second would give.
It holds on residuals and at $d = 10$, but not at $d = 3$. The drop persists after controlling for block information: the gain over a prior-only model falls from 0.074 to 0.039.}
\end{frame}

% =====================================================================================
\begin{frame}{Content (P2): the prior does not rise with horizon}
\begin{columns}[T,onlytextwidth]
\column{0.55\textwidth}
\centering
\begin{tikzpicture}
\begin{axis}[subspace bars, ymin=-0.2, ymax=0.12, ytick={-0.2,-0.1,0,0.1}, ylabel={prior, corrected \Rsq{}}]
\draw[cInk, line width=0.6pt] (rel axis cs:0,0.625) -- (rel axis cs:1,0.625);
% results/stats/d5/matrix.csv, variable prior (mean and 95% CI over 13 animals)
\addplot[fill=cRRR, draw=none, bar shift=0pt, ci bars] coordinates {
  (r1,0.054) += (0,0.051) -= (0,0.048)
  (r4,-0.027) += (0,0.086) -= (0,0.092)
  (r16,-0.057) += (0,0.112) -= (0,0.122)
  (r60,-0.046) += (0,0.107) -= (0,0.110)};
\addplot[fill=cGrey, draw=none, bar shift=0pt, ci bars] coordinates {
  (pca,0.013) += (0,0.033) -= (0,0.040)
  (rnd,-0.030) += (0,0.030) -= (0,0.031)};
\end{axis}
\end{tikzpicture}

\muted{\small Corrected = \Rsq{} minus pseudo-session null median; 95\% CI over mice}

\column{0.41\textwidth}
Change from 0.5 s to 30 s:\\
\hl{$-0.099$} [$-0.218$, $+0.002$]\\
Holm $p = 0.89$; 4 of 13 mice rise.

\medskip
Weak at every lag: no lag significant across sessions (best Fisher $p = 0.12$).

\medskip
No rise on residuals (there significantly the \emph{wrong} way), at $d = 3$ or 10, or in any leave-one-mouse-out
refit.
\end{columns}
\note{%
This is the result I find most interesting. Same layout, now for the Bayes-optimal prior, where the hypothesis
predicted bars growing from left to right. They do not: the change from short to long goes the wrong way, in every
variant I ran. Prior decoding from the subspaces is weak at every lag, while the full population carries it. So the
prior is not concentrated in the predictable directions. Why can bars go below zero? Decoding uses contiguous
cross-validation folds, which penalise ramp-dominated features, so the real decoder can score below the null median.
\qa{The 0.5 s bar is above zero; isn't the prior there?}{Weakly: its CI [0.006, 0.105] excludes 0, but the session-level
combined test is not significant ($p = 0.15$). Either way it is the short lag, opposite to the hypothesis.}
\qa{Could the prior be in the 30 s subspace after all?}{Its CI [$-0.156$, $+0.061$] is wide, so a full-size effect
there is not excluded. What is clear is that it does not rise with horizon.}}
\end{frame}

% =====================================================================================
\begin{frame}{What the long horizon finds instead: a slow mode that tracks time in session}
\centering
\figorplaceholder{../figures/slide_fig5.pdf}{0.98\textwidth}{5.2cm}

\raggedright\small
Mostly a ramp (median $r$ 0.96), sometimes a step. Time \Rsq{}: 30 s \hl{0.84}, 0.5 s 0.82, PCA 0.80, random 0.64.
\hl{Drift flag} (time above PCA and random, prior n.s.): \hl{14 of 16}. Nested: \hl{92\%} inside the 0.5 s subspace.
\note{%
What does the 30 s subspace carry? Time. Its leading direction, the slow mode, in 30 s windows: the session with the
median correlation with time, and two step-like sessions that look more like discrete events. Elapsed-time \Rsq{} is
0.84, but PCA gets 0.80 and random subspaces 0.64: every low-dimensional summary here is dominated by the session-scale
ramp, long horizons only slightly more. The 0.5 s subspace contains 92\% of this direction.
\qa{How is elapsed time decoded?}{Ridge on the latents; contiguous chunks assigned round-robin to folds, so it
interpolates; subspaces are fit on the whole session. Hence the PCA and random comparison.}
\qa{Isn't this just a failed hypothesis?}{The proposal planned for it: ``slow subspace is drift'' $\to$ report a
negative result on long-lag objectives; ``no dissociation across lags'' $\to$ report the shared slow mode and
characterise it. The pilots do that.}}
\end{frame}

% =====================================================================================
\begin{frame}{Is it just recording drift? Most controls barely dent the slow mode}
\centering
\begin{tikzpicture}
\begin{axis}[slide axis, width=10.0cm, height=5.0cm, xmin=0, xmax=1, ymin=0.4, ymax=6.6, y dir=reverse,
  xtick={0,0.2,0.4,0.6,0.8,1}, ytick={1,...,6}, ymajorgrids=false, xmajorgrids, enlarge x limits=0, clip=false,
  yticklabels={{Original units (16 / 13)},{Rate-stable units (13 / 12)},{Multi-unit, no sorting (15 / 12)},
    {Amplitude-stable units (5 / 5)},{Amplitude-corrected, maximal},{Slow-maximal, post hoc}},
  xlabel={elapsed-time \Rsq{} of the 30 s subspace},
  error bars/x dir=both, error bars/x explicit, error bars/error bar style={line width=1.4pt}]
% thresholds: 75% and 50% of the original 0.8413 (pre-registered "survives" / "removed")
\draw[cInk, dashed] (axis cs:0.6310,0.4) -- (axis cs:0.6310,6.6);
\draw[cInk, dashed] (axis cs:0.4206,0.4) -- (axis cs:0.4206,6.6);
\node[font=\small, color=cInk, anchor=south] at (axis cs:0.6310,0.4) {75\%};
\node[font=\small, color=cInk, anchor=south] at (axis cs:0.4206,0.4) {50\%};
% same-session original (hollow), from results/stats/d5/drift.csv restricted to each version's sessions
\addplot[only marks, mark=o, mark size=3pt, cGrey, line width=1pt] coordinates {
  (0.8413,1) (0.8340,2) (0.8790,3) (0.7339,4) (0.8413,5) (0.8413,6)};
% control means and 95% CI over mice, results_*/stats/d5/matrix.csv (rrr_60, elapsed_time)
\addplot[only marks, mark=*, mark size=3pt, cRRR, error bars/error bar style={cRRR, line width=1.4pt}] coordinates {
  (0.8413,1) += (0.0650,0) -= (0.0908,0)
  (0.7735,2) += (0.0723,0) -= (0.0907,0)
  (0.7424,3) += (0.0843,0) -= (0.1125,0)
  (0.4933,4) += (0.1510,0) -= (0.1056,0)
  (0.7959,5) += (0.0738,0) -= (0.0954,0)};
\addplot[only marks, mark=*, mark size=3pt, cPCA, error bars/error bar style={cPCA, line width=1.4pt}] coordinates {
  (0.7826,6) += (0.0799,0) -= (0.1042,0)};
\end{axis}
\end{tikzpicture}

\vspace{-1mm}
\raggedright\small Filled: control, 95\% CI; hollow: original, same sessions; dashed: 75\% / 50\% of 0.84.\\
Of the same sessions' original: rate-stable 93\%, multi-unit 84\%, \hl{amplitude-stable 67\%} (PCA there: 0.43),
amplitude-corrected 93--95\%.
\note{%
The obvious worry is electrode drift: the probe moves, units fade, spikes shrink, and firing seems to ramp. Each row is
one version of the data; filled dots are the control, hollow ones the original on the same sessions, and the dashed
lines are the pre-registered 75\% ``survives'' and 50\% ``removed'' thresholds. Most controls barely move it. The
exception is the strictest filter, units stable in both rate and amplitude: only 5 sessions pass, it is lower in all 5,
and it loses a third of its value against the same sessions, two-fifths against the full original. PCA also falls
there, to 0.43, so part of the drop is simply having fewer units.
\qa{So it is not drift?}{No single recording cause I tested accounts for most of it. Untested ones remain possible,
such as settling before the task started.}
\qa{Where is the conservative amplitude correction?}{Omitted here: its coefficient is negative ($-1.43$), so it corrects
in the wrong direction for detection loss and does not test it (backup).}}
\end{frame}

% =====================================================================================
\begin{frame}{Candidate explanations, tested one at a time}
\small
\setlength{\tabcolsep}{5pt}
\begin{tabular}{@{}L{0.29\textwidth}L{0.33\textwidth}L{0.31\textwidth}@{}}
\toprule
\textbf{Candidate} & \textbf{Test} & \textbf{Result} \\
\midrule
Trial-locked responses & refit on GLM residuals & \Rsq{} 0.73 (86\%) \\
Units entering, leaving & rate-stable units only & \Rsq{} 0.77 (93\%) \\
Spike-sorting errors & multi-unit activity, no sorting & \Rsq{} 0.74 (84\%); same mode \\
Amplitude, detection loss & coupling; corrected refits & wrong sign; 85--86\% of ramp left \\
\bottomrule
\end{tabular}

\medskip
\normalsize No single recording cause tested accounts for most of it; the strictest unit filter (5 sessions) cuts it by
a third to two-fifths.

\medskip
\muted{\small Also tested (backup): probe motion, pupil, GLM-HMM engagement, satiety; none tied to the mode beyond the
trend. \% = of the same sessions' original. Pilots: exploratory, after the analysis freeze, each pre-registered.}
\note{%
The scorecard uses several tools from this course: a GLM to remove trial-locked responses, sorting-free multi-unit
activity to bypass spike sorting, and a GLM-HMM engagement state. Amplitude was the interesting one: units whose spikes
shrank were the ones whose rate fell, which looked like spikes slipping under the detection threshold. Tested directly,
that failed: the amplitude--rate coupling has the wrong sign, and even a maximal correction leaves 85 to 86\% of the
ramp. The A13 rule said ``sorting artifact''; the multi-unit result contradicts that label, and A14 refuted the
detection-loss reading. Each test could have removed the mode; none did.
\qa{Why not call it arousal or fatigue?}{The mode is mostly a near-straight ramp, so any slowly trending signal
matches it about equally well; these tests have little power to tell causes apart.}}
\end{frame}

% =====================================================================================
\begin{frame}{Does it generalise? Other areas on the same probes}
\begin{columns}[T,onlytextwidth]
\column{0.54\textwidth}
\small
% results/a15/summary_d5.json; results_region2/stats/d5/{primary,secondary,drift}.csv; MOs from results/stats/d5
\setlength{\tabcolsep}{4pt}
\begin{tabular}{@{}lcc@{}}
\toprule
 & MOs & Other area \\
\midrule
Sessions, mice & 16, 13 & 15, 12 \\
P1 geometry, Holm $p$ \hl{*} & 0.0026 & 0.0037 \\
P2 prior rises, Holm $p$ \hl{*} & 0.89 & 0.52 \\
P3 stimulus falls, Holm $p$ \hl{*} & 0.0061 & 0.012 \\
Elapsed-time \Rsq{}, 30 s \hl{*} & 0.84 & 0.82 \\
Drift flag & 14 / 16 & 14 / 15 \\
Nested in 0.5 s subspace & 0.92 & 0.92 \\
\midrule
Prior, population \Rsq{} & 0.064 & 0.001 \\
\bottomrule
\end{tabular}

\smallskip
\hl{*} Pre-registered: P1, P3 significant, P2 not, \Rsq{} $\ge 0.5$. \hl{All four hold.}

\column{0.42\textwidth}
\small
Largest non-MOs area per session: frontal (ACAd, ILA, PL, ORBl, ORBvl, AIv, MOp) in 10 sessions; septum and striatum in
5.

\smallskip
``P2 not significant'' says little here: the population prior is about 0.

\smallskip
MOs and other-area slow modes share no fluctuations beyond the ramp ($p = 0.48$).

\smallskip
The MOs mode (16 sessions, 13 mice) tracks neither reaction time ($p = 0.12$) nor accuracy ($p = 0.34$) beyond the trend.
\end{columns}
\note{%
Is any of this specific to MOs? The probes pass through other areas, so in each session I took the largest one and
reran the unchanged pipeline, with criteria registered before binning. All four hold: geometry and stimulus effects
replicate, and the same slow mode dominates, nested in the short-lag subspace. Two caveats. The prior is nearly absent in
these areas, so ``no rise'' there is uninformative. And after removing the ramp, the two areas' slow modes share no
fluctuations, and the MOs mode tracks neither reaction time nor accuracy. Still, the pipeline transfers unchanged, which
is the evidence that scaling up is worthwhile.
\qa{Same probe: isn't this the same artifact twice?}{Yes, the ramp itself could be a probe-wide effect; this replication
shows the pattern is not specific to MOs, not that it is neural. But the fluctuations around the ramp are not shared
even on the same probe (median detrended $r$ 0.07): weak evidence against a probe-level cause.}}
\end{frame}

% =====================================================================================
\begin{frame}{What this could mean for how cortex carries beliefs and state}
\begin{columns}[T,onlytextwidth]
\column{0.48\textwidth}
\textbf{The prior: present, not predictive}

\medskip
MOs holds the block belief, but it is \hl{not concentrated in the predictable directions} and does not rise with
horizon.

\medskip
One reading: a \hl{small, static offset} updated by trial events, not a slow mode the dynamics carry forward.

\medskip
\muted{Variance- or predictability-ranked methods would then miss it.}

\column{0.48\textwidth}
\textbf{The slow mode: a session-long change}

\medskip
Seen in MOs and in other areas on the same probes; no single recording cause tested accounts for most of it.

\medskip
Candidates: arousal or neuromodulation, fatigue, satiety, elapsed-time coding, or an untested probe-wide factor.

\medskip
\muted{Either way, slow cognitive variables must be separated from it.}
\end{columns}
\note{%
What could this mean for neuroscience? These are interpretations, not results. First, MOs holds the block belief, but
not in the directions that predict its own future. One reading is a small, static offset updated by trial events,
rather than a slow mode the circuit keeps alive; if so, methods that rank dimensions by variance or predictability will
miss it. Second, the slow mode is a session-long change that no single recording cause I tested explains. Arousal,
neuromodulation, fatigue or a sense of elapsed time are candidates, as is a probe-wide factor I did not test. Either
way, it has to be separated from any slow cognitive variable.}
\end{frame}

% =====================================================================================
\section{Conclusion}
\begin{frame}{What I can and cannot conclude}
\begin{columns}[T,onlytextwidth]
\column{0.48\textwidth}
\textbf{Supported}
\begin{itemize}\setlength{\itemsep}{3pt}
\item Prior decoding does not rise with horizon; it is weak in every subspace.
\item Predictable 2--30 s structure is one slow mode, shared across horizons, tracking time in session.
\item No single recording cause tested accounts for most of it.
\item The pattern repeats in other areas on the same probes.
\end{itemize}

\column{0.48\textwidth}
\textbf{Open, and limitations}
\begin{itemize}\setlength{\itemsep}{3pt}
\item What the mode \emph{is}; the strictest unit filter (5 sessions) cuts it by a third to two-fifths.
\item 13 mice; other areas share probes; pilots exploratory, uncorrected.
\item Split halves confound reliability with nonstationarity.
\item A ramp limits every behavioural comparison.
\end{itemize}
\end{columns}
\note{%
To be clear about the claims: prior decoding does not rise with horizon; the predictable long-horizon structure is a
single slow mode that tracks time; no single recording cause I tested explains most of it; and the pattern repeats in
other areas. What the mode is remains open. The strictest unit filter does weaken it, and because the other areas share probes with
MOs, a probe-wide factor I did not measure is still possible. This is 13 mice, and the pilots are exploratory and
uncorrected, so I treat them as characterisation rather than confirmation.
\qa{What would change your mind?}{A recording measure that follows the mode's session-specific shape, steps included,
would point to an artifact; the mode continuing through task-free periods would point away from task-related state.}}
\end{frame}

% =====================================================================================
\begin{frame}{Takeaways and next steps}
\begin{enumerate}\setlength{\itemsep}{2pt}
\item The block prior is decodable from MOs, but it \hl{does not rise with prediction horizon}; subspace prior decoding
is weak at every lag (none significant across sessions).
\item Long horizons find \hl{one slow mode that tracks time in session}, shared with the short horizon.
\item \hl{No single recording cause tested accounts for most of it}; its origin is open.
\item The pattern \hl{repeats in other areas on the same probes} (mostly frontal).
\end{enumerate}

\textbf{Next.} Does the mode continue through the task-free spontaneous period? Then independent sessions and other
regions (MOp: 14 sessions, 12 mice).

\muted{\small Lesson: long-horizon objectives need drift controls (elapsed-time decoding, pseudo-session nulls).}
\note{%
To sum up: the prior does not rise with horizon; long horizons find one slow, time-tracking mode; no single recording
cause I tested explains most of it; and the pattern repeats in other areas on the same probes. The most informative
next test is the spontaneous period after the task: if the mode keeps ramping when the mouse is not doing the task, it
is not tied to the task. Then independent sessions and other regions. The practical lesson: a long-horizon predictive
objective is a very good drift detector, so it needs drift controls. Thanks, happy to take questions.
\qa{What would you do with more time?}{Register and run the spontaneous-epoch test; download independent sessions for
MOp and other regions; try supervised, null-controlled methods aimed at small static codes.}}
\end{frame}

% =====================================================================================
\appendix
\begin{frame}[standout]
Backup slides
\note{Backup slides for questions: data and statistics, why RRR, the prior result, robustness, the full decoding matrix,
nesting, the amplitude pilot, probe motion and state signals, and the normalisation bug.}
\end{frame}

% -------------------------------------------------------------------------------------
\begin{frame}{Backup: data and statistics}
\small
\begin{columns}[T,onlytextwidth]
\column{0.48\textwidth}
\textbf{Region.} MOs was the only candidate meeting the frozen rule ($\ge 15$ sessions, $\ge 8$ mice): MOs 16/13, MOp
14/12, VISp 10/8, ACAd 4/4. MOp's published prior effect is slightly larger.

\medskip
\textbf{Data.} 16 sessions, 13 mice. 22--128 good units per session (median 48.5), at least 0.1 Hz.
413--700 trials (median 534) after the trial mask. Task epoch only.

\medskip
\textbf{Fits.} 0.5 s bins of $\sqrt{\text{counts}}$, z-scored. Contiguous cross-validation folds, purged and embargoed
around test blocks.

\column{0.48\textwidth}
\textbf{Null.} 200 pseudo-sessions, shared across lags and feature sets, so lag contrasts get a paired null. Corrected
score = real minus null median.

\medskip
\textbf{Tests.} Sessions averaged within mouse; one-sided Wilcoxon signed-rank across 13 mice; Holm over the 3
primary tests; geometry as an intersection--union test. Smallest attainable $p$ with 13 mice: $1.2\times10^{-4}$.

\medskip
\textbf{Pilots.} A12--A15: exploratory; A13--A15 ran after the 20 Nov analysis freeze. Each pre-registered; no
multiplicity correction.
\end{columns}
\note{%
The unit of inference is the mouse: 16 sessions from 13 mice, averaged within mouse before every test.
\qa{Why one-sided tests?}{The proposal said two-sided; I switched to one-sided in the pre-data amendment (D6) because
the hypotheses are directional. Two-sided doubles the $p$-values: Holm P1 0.0051, P3 0.012, still significant.}
\qa{How is the code organised?}{A package, src/pns (binning, lagged RRR, decoding, geometry, prior, slow mode), driven
by numbered scripts 01 to 20 in scripts/, run in order from the census to the figures and pilots, with unit tests in
tests/. Pre-registrations and deviations are in AMENDMENTS.md, runs in ANALYSIS\_LOG.md.}
{\footnotesize\textbf{Code-level prep.} Folds: the subspace decoders use contiguous folds (blocks of trials or time,
purged around test blocks); the full-population gate uses Findling's interleaved (shuffled) folds with an inner CV over
the Lasso penalty. Contiguous folds are stricter for slow signals: a ramp learned on one block extrapolates badly to
another. Drift flag (A10): 30 s elapsed-time \Rsq{} above PCA and random, and 30 s prior not significant.\par}}
\end{frame}

% -------------------------------------------------------------------------------------
\begin{frame}{Backup: why lagged RRR rather than PCA or other methods?}
\small
\begin{itemize}\setlength{\itemsep}{1pt}
\item \textbf{The horizon is an explicit knob.} One objective; between fits only $\Delta$ changes (same rank, same
ridge penalty).
\item \textbf{Simple and well studied.} Ridge fit plus an SVD (Izenman 1975; Semedo et al.\ 2019).
\item \textbf{CCA.} With a ridge penalty, RRR sits between CCA-like (small $\lambda$) and PLS-like (large $\lambda$)
solutions, which is why $\lambda$ was fixed across lags.
\item \textbf{PCA} ranks by variance and has no notion of time. \textbf{DCA} (Clark et al.\ 2019) ranks by
predictability over a window, not at a chosen horizon. Time-lagged ICA and slow features are close cousins.
\item \textbf{GPFA} learns per-latent timescales in a generative model rather than optimising prediction at a chosen
horizon (a proposal extension not reached).
\item \textbf{Empirically.} RRR beats PCA at every lag, but by only 0.011--0.022 in \Rsq{}; on trial-locked residuals
the margin is 0.005--0.008 and its CIs include 0.
\end{itemize}
\note{%
I chose lagged reduced-rank regression because it makes the horizon an explicit parameter while everything else stays
fixed, and it is closed form. The honest empirical answer is that it beats PCA only modestly, and much of that edge is
task structure that recurs every trial, since it nearly vanishes on residuals. That fits the main finding: most
predictable structure here is either trial-locked or the one slow mode, which PCA also finds.}
\end{frame}

% -------------------------------------------------------------------------------------
\begin{frame}{Backup: how strong is the prior result?}
\small
\begin{itemize}\setlength{\itemsep}{4pt}
\item \textbf{Positive controls pass.} The prior is decodable from the full population (Fisher $p = 1.2\times10^{-5}$),
and stimulus decoding and its fall with lag are detected (P3).
\item \textbf{The pre-registered contrast.} Change from 0.5 s to 30 s: $-0.099$ [$-0.218$, $+0.002$] (\Rsq{} units).
The population gate uses a different decoder (Lasso on all units, interleaved folds), so no ratio is taken.
\item \textbf{No rise in any variant.} Residuals: significantly the wrong way, $-0.093$ [$-0.154$, $-0.038$].
$d = 3$ and 10: Holm $p = 0.21$, 0.88. Leave-one-mouse-out: all Holm $p \ge 0.80$. Pilot refits: 0.63--1.0. Other
areas: 0.52.
\item \textbf{Limit.} Per-lag CIs are wide: at 30 s [$-0.156$, $+0.061$], so a full-size prior there is not excluded.
The claim is about the change with horizon.
\end{itemize}
\note{%
The pipeline detects things when they are there: the prior from the full population, and the stimulus in the
subspaces. The pre-registered contrast never shows the predicted rise in any variant. The limit is precision at each
lag: the 30 s interval is wide.
{\footnotesize\textbf{Code-level prep.} paired\_null\_diff (src/pns/analysis.py): P2 and P3 test, per session, the
difference between the two lags' raw scores, corrected by the median of the same difference across pseudo-sessions,
not the difference of two separately corrected scores. The null is shared across lags, so the pairing cancels
session-specific offsets.\par}}
\end{frame}

% -------------------------------------------------------------------------------------
\begin{frame}{Backup: robustness of the three primary tests (Holm $p$)}
\centering
\small
\begin{tabular}{@{}lccc@{}}
\toprule
Variant & P1 geometry & P2 prior & P3 current side \\
\midrule
Primary analysis ($d = 5$, 13 mice) & 0.0026 & 0.89 & 0.0061 \\
Trial-locked residuals & 0.013 & 0.996 & 0.0051 \\
Leave-one-mouse-out, worst case (raw) & 0.005 & 0.95 & 0.012 \\
Leave-one-mouse-out, worst case (residual) & 0.027 & 0.998 & 0.010 \\
Drop outlier mouse DY\_014 & 0.005 & 0.87 & 0.005 \\
$d = 3$ & 0.032 & 0.21 & 0.11 \\
$d = 10$ (10 mice) & 0.003 & 0.88 & 0.006 \\
Rate-stable units only (12 mice) & 0.002 & 0.97 & 0.027 \\
Maximal amplitude correction & 0.0018 & 0.71 & 0.013 \\
Other areas on the same probes (12 mice) & 0.0037 & 0.52 & 0.012 \\
\bottomrule
\end{tabular}

\muted{\small Residuals: planned add-on. Leave-one-out, drop DY\_014: post hoc. Last three rows: pilots.}
\note{%
Every row keeps the same pattern, geometry and current side significant and prior not, except current side at
$d = 3$. DY\_014 is the low outlier in the geometry analysis: one 141 Hz cluster with very high autocorrelation that
loses about 40\% of its rate late in the session, likely a multi-unit or noise cluster plus a drift event. Its effects
point the same way, and dropping it changes nothing material.
\qa{Why only 10 mice at $d = 10$?}{Sessions need enough units to fit 10 dimensions, so 12 sessions from 10 mice
remain.}}
\end{frame}

% -------------------------------------------------------------------------------------
\begin{frame}{Backup: full decoding matrix ($d = 5$, mean over 13 mice)}
\centering
\small
% results/stats/d5/matrix.csv (feature sets rrr_1, rrr_4, rrr_16, rrr_60, pca, random); 95% CIs in figures/fig2_content_d5.pdf
\setlength{\tabcolsep}{4pt}
\begin{tabular}{@{}lcccccc@{}}
\toprule
 & RRR 0.5 s & RRR 2 s & RRR 8 s & RRR 30 s & PCA & random \\
\midrule
Current side (AUC, corrected) & 0.218 & 0.161 & 0.144 & 0.158 & 0.240 & 0.218 \\
Previous side (AUC, corrected) & 0.045 & 0.050 & 0.044 & 0.037 & 0.027 & 0.007 \\
Prior (\Rsq{}, corrected) & 0.054 & $-0.027$ & $-0.057$ & $-0.046$ & 0.013 & $-0.030$ \\
Elapsed time (\Rsq{}) & 0.821 & 0.834 & 0.839 & 0.841 & 0.800 & 0.641 \\
\midrule
Choice (AUC) & 0.615 & 0.604 & 0.589 & 0.590 & 0.638 & 0.616 \\
Feedback (AUC) & 0.803 & 0.708 & 0.600 & 0.609 & 0.807 & 0.790 \\
Wheel speed (\Rsq{}) & 0.339 & 0.191 & 0.024 & 0.043 & 0.406 & 0.272 \\
\bottomrule
\end{tabular}

\medskip
\muted{Corrected = real minus pseudo-session null median. Lower rows depend on the animal, so no null applies.}
\note{%
The whole variables-by-subspaces matrix. The upper four rows are the core ladder, ordered by timescale; the prior row
is the one that should have turned positive at long lags and did not. Choice, feedback and wheel speed depend on the
animal, so the pseudo-session null does not apply to them; they are raw scores, read descriptively. Feedback and wheel
speed fade with lag like the stimulus does.}
\end{frame}

% -------------------------------------------------------------------------------------
\begin{frame}{Backup: one shared slow direction (nesting and projected subspaces)}
\centering
\figorplaceholder{../figures/slide_figS1.pdf}{\textwidth}{4.8cm}

\raggedright\small
(a) The 0.5 s subspace contains 92\% of the 30 s leading direction; the reverse is 76\%. Chance here is isotropic, a
low bar; no covariance-matched chance was computed. (b) ``30 s off 0.5 s'' is the 30 s subspace projected off the
0.5 s subspace: a new 5-D space, not part of the 30 s one, so it can score higher (stimulus 0.220 vs 0.158).
\note{%
The nesting index is the squared length of the 30 s leading direction after projection onto the 0.5 s subspace, so 1
means fully inside. The slow direction is shared rather than split off. Against isotropic chance that is a low bar;
PCA's leading direction would likely also sit inside, so read it as ``the same slow direction'', not as evidence of
something special.
{\footnotesize\textbf{Code-level prep.} complement\_basis(A, B) (src/pns/geometry.py) orthonormalises the part of
span(A) orthogonal to span(B): $Q_A - Q_B Q_B^\top Q_A$, then an SVD. The result is generally not inside span(A).
Decoding scores also depend on the basis, because features are standardised and L2-penalised, so a projected space
can decode a variable better than the subspace it came from.\par}}
\end{frame}

% -------------------------------------------------------------------------------------
\begin{frame}{Backup: does amplitude loss account for the slow mode? (third pilot)}
\begin{columns}[T,onlytextwidth]
\column{0.44\textwidth}
\small
\textbf{Coupling beyond the trend.} Per-unit Poisson GLM: 30 s-window count on relative spike amplitude, quadratic
time as nuisance. Detection loss predicts a positive coefficient.

\medskip
Observed: mean \hl{$-1.54$}; 27\% of units positive; $p = 0.99$.
Isolated spikes only (post hoc): $-1.32$, $p = 0.98$.

\medskip
\textbf{Timing.} Settling predicts early change: index $-0.12$ ($p = 0.75$).

\column{0.52\textwidth}
\small
\setlength{\tabcolsep}{4pt}
\begin{tabular}{@{}lcc@{}}
\toprule
Correction & time \Rsq{} & ramp left \\
\midrule
Conservative$^\dagger$ & 0.85 (101\%) & 1.00 \\
Maximal & 0.80 (95\%) & 0.86 \\
Slow-max.\ (post hoc) & 0.78 (93\%) & 0.85 \\
\bottomrule
\end{tabular}

\medskip
$^\dagger$Its coefficient is negative ($-1.43$), so it corrects the wrong way: not a test. Primary tests unchanged in
every copy.

\medskip
\muted{Caveats: factors clipped to 0.5--2; Kilosort amplitudes are template scalings; settling before task
start is invisible.}
\end{columns}
\note{%
The second pilot found that units whose spikes shrank were the units whose rate fell along the mode, which looked like
detection loss. This pilot tested sufficiency. The coupling, estimated only from fluctuations around the trend, is
negative: in a window where a unit's spikes are 10\% smaller, it fires about 16\% more, the opposite of detection loss.
It stays negative using only isolated spikes, so bursts are not hiding it. Removing from every unit everything that
moves with its own amplitude still leaves 93 to 95\% of the elapsed-time decoding.}
\end{frame}

% -------------------------------------------------------------------------------------
\begin{frame}{Backup: probe motion, sorting and state signals (first two pilots)}
\small
\begin{columns}[T,onlytextwidth]
\column{0.48\textwidth}
\textbf{Probe motion} (IBL estimate; median range 7.9\,\textmu m). $|r|$ with own trace 0.74, but no higher than with
other sessions' traces: correction $-0.006$, $p = 0.39$; detrended $+0.062$, $p = 0.14$.

\medskip
\textbf{Rigid-shift signature.} Multi-unit loadings vs.\ depth-profile slope: mean $|\rho| = 0.19$, combined
$p = 0.25$.

\medskip
\textbf{Sorting.} Multi-unit mode matches the unit mode (median $r = 0.96$; 15 sessions).

\column{0.48\textwidth}
\textbf{Pupil} (15 sessions): $|r| = 0.35$; correction $-0.001$, $p = 0.42$; detrended $p = 0.55$.

\medskip
\textbf{Engagement} (GLM-HMM), beyond a quadratic trend: $p = 0.19$; fluctuating subset $p = 0.049$
(uncorrected), sign inconsistent ($p = 0.95$).

\medskip
\textbf{Satiety} (cumulative rewards): $p = 0.12$; sign leans positive ($p = 0.068$).

\medskip
\textbf{Amplitude--loading link:} real but small (Fisher $z$ 0.15, $p = 0.029$).
\end{columns}
\note{%
The logic of the session-permutation tests: correlate a session's slow mode with every other session's signal, and ask
whether its own signal fits better than other sessions' do. For probe motion and pupil it does not. The catch is that
two ramps always correlate, so these tests only separate causes that produce different shapes; that is why the later
pilots moved to evidence that does not depend on matching ramps.}
\end{frame}

% -------------------------------------------------------------------------------------
\begin{frame}{Backup: the normalisation bug (amendment D10)}
\small
\begin{itemize}\setlength{\itemsep}{5pt}
\item \textbf{Found} while diagnosing two sessions where every method, full-rank ridge included, had held-out
\Rsq{} near 0.
\item \textbf{Cause.} The z-score SD floor ($10^{-3}$) let a unit silent in a training fold but active in its test fold
reach a held-out z-variance of about $10^{5}$, over 99\% of pooled test variance. That corrupted cross-validated
\Rsq{} and, through it, the session ridge penalty.
\item \textbf{Fix.} Floor = SD of a unit firing exactly at the 0.1 Hz inclusion threshold (0.218). No whole-session
z-score changes (0 of 842 units floored).
\item \textbf{Effect.} Ridge penalty changed in 2 sessions (raw) and 4 (residual). Raw primary $p$-values unchanged.
P1 moved: residual 0.048 to 0.013; residual leave-one-mouse-out worst case 0.092 (n.s.) to 0.027; $d = 3$ 0.049 to
0.032. ``RRR beats full-rank ridge at 30 s'' was an artifact and was dropped.
\end{itemize}
\note{%
This is the one analysis decision made after seeing results, and it is a bug fix. It was exposed by unit turnover:
units that go silent for a whole fold. The subspaces themselves use whole-session z-scoring and were unaffected; the
damage was to fold-wise scoring and the penalty choice. All pre-fix outputs are kept in the repository for comparison.
\par\bigskip\textbf{Reminder:} if asked, say that AI tools were used in this project and are disclosed in the report.}
\end{frame}

\end{document}
